% !PRINTED 767
\weberpage{767}
But by 5. and 6. it has been shown that, if $p$ is large enough, equation (15) cannot hold. Thereby 4., and with it the whole of Theorem I., is proved.

This theorem now admits manifold applications. It gives us first of all the transcendence of $e$, if we take $C_1, C_2, \ldots, z_1, z_2, \ldots$ to be integral rational numbers. It further gives the transcendence of $\pi$. For from the equation $1+e^{i\pi}=0$ it follows by I. that $i\pi$, and consequently $\pi$ also, cannot be algebraic.

From this there follows further:

\gesperrt{For every algebraic number $x$, with the exception of $x=0$, $X=e^x$ is a transcendental number.}

\gesperrt{For every algebraic $X$, with the exception of $X=1$, every natural logarithm $x=\log X$ is a transcendental number.}

\gesperrt{For every arc that stands to the radius in a ratio $x$ expressible algebraically, with the exception of $x=0$, $X=\sin x$ is a transcendental number.}

This follows by I. from $2i\,X=e^{ix}-e^{-ix}$.

% [sic] -- the source prints the chord as $\frac12\sin(x/2)$; for unit radius the chord is $2\sin(x/2)$.
% [sic] -- the source's claim for $\mathrm{tg}\,x=\alpha x$ omits the algebraic root $x=0$; the intended claim concerns nonzero roots.
The same holds for the other trigonometric lines, $\cos x$, $\mathrm{tg}\,x$, and for the chord $\dfrac{1}{2}\sin\dfrac{x}{2}$. To adduce one thing more: the transcendental equation $\mathrm{tg}\,x=\alpha x$ has, for an algebraic $\alpha$ other than 0, only transcendental numbers for roots.

\begin{center}
\rule{4cm}{0.4pt}
\end{center}
